\setlength{\parindent}{0pt}
% ------------------------------------------------------------
\chapter{Future Work}
\label{chap:futurework}
% ------------------------------------------------------------
The tool as it stands leaves several research directions unaddressed. Some of these directions
sharpen the engine
without changing the model it verifies, others widen the model, making the hardware
it protects represent a larger, more realistic class. We collect them here, from the most immediate to the most
ambitious.

\section{Arithmetic Circuits on Larger Fields}
\label{sec:fw-fields}

It is tempting to file work over larger fields under ``not yet supported,'' but
much of the machinery is already field-agnostic and it is worth saying exactly
is supported and what is not. The coupling characterisation of Chapter~\ref{chap:couplings} and the
subset-space traversal of Section~\ref{sec:subset-space} never used any property of
$\mathbb{F}_2$ beyond a finite randomness space sampled uniformly, so they carry
over to any finite field trivially. The substitution at the heart of the engine
generalises just as cleanly. Over $\mathbb{F}_q$ the invertible affine map
$r \mapsto c \cdot r + e$ with $c \neq 0$ and $r \notin \mathrm{vars}(e)$ is
also measure-preserving on the randomness space and plays exactly the role the shear
$r \mapsto e \oplus r$ plays over $\mathbb{F}_2$.

What is genuinely tied to $\mathbb{F}_2$ is the concrete engine built on that
substitution. The exact endpoint test when extending a probe set rests on the algebraic normal form, whose
multilinearity is the $\mathbb{F}_2$ identity $x^2 = x$; the multiplicative
factoring of Section~\ref{sec:factoring} and the canonicalisation it runs through
are likewise $\mathbb{F}_2$ algebra, as is the two-input XOR/AND gate basis. So the
step to a larger field necessitates to build on top rather than to redesign. One can replace
the endpoint with a suitable normal form, or with a Schwartz--Zippel probabilistic
identity test (though that would harm the exactness), which is in fact cheaper in the larger field,
re-derive factoring
for the new algebra, and enlarge the gate set. Verification is already moving this
way, toward arithmetic and prime-field masking under the pressure of post-quantum
schemes \parencite{gigerl2023formal} which work on larger fields, and a field-generic version of our coupling
engine would fit naturally into that line.

\section{Factoring Heuristics}
\label{sec:fw-factoring}

The multiplicative factoring of Section~\ref{sec:factoring} pulls out, greedily,
the factor shared by the most product terms, breaking ties by identifiers.
Factoring is not confluent: which factor one extracts, and in what order relative
to the substitutions, decides whether a buried random surfaces and how much
cancellation follows. The present choice is cheap and suffices on our benchmarks,
but it is surely not optimal. One might weight a candidate factor by whether
pulling it out actually exposes an additive single-occurrence random rather than
merely by how often it appears, i.e. it can explore a shallow lookahead over factor choices, or
infer an order from the gadget's structure. Because factoring preserves a
wire's value (Section~\ref{sec:enlarging-verifiable}), every such heuristic is
sound by construction; the question is purely how often, and how quickly, it
reaches a discharge. A systematic study of these heuristics, measured against the
completeness they buy, is left open.

\section{Composing Couplings for Unions of Safe Sets}
\label{sec:fw-coupling-union}

Component~III certifies one safe set from one coupling. When the traversal
certifies probes on sibling branches it obtains two safe sets $S_1$ and $S_2$, each
internally secret-independent, with couplings $\varphi_1$ and $\varphi_2$. Their
union need not be secret-independent of course, and the present scheme makes no attempt on
it; the branches are kept apart. The couplings themselves, though, may sometimes be
composed into a certificate for the union.

Call the \emph{randomness footprint} of a safe set the randoms its wires depend on
together with the randoms its coupling relabels. Suppose $S_1$ and $S_2$ have
disjoint footprints. Then $\varphi_1$ and $\varphi_2$ move disjoint coordinates of
the randomness space, so they commute nicely, and neither disturbs the wires the other has
already blinded. Their composite $\varphi_1 \circ \varphi_2$ is again
measure-preserving, and it carries every wire of $S_1 \cup S_2$ to a secret-free
form. In such a case, a wire of $S_1$ is blinded by $\varphi_1$ and left untouched by $\varphi_2$,
and symmetrically for $S_2$. By the safe set fact of
Section~\ref{sec:coupling-extension} that single composite certifies the whole
union, merging two safe sets with no fresh discharge.

The disjointness is the key here, and it is why the idea is a direction and not
yet a method. When the footprints overlap, as two gadgets sharing a mask, the
shears no longer commute and one coupling can undo the blinding the other needs,
so the composite is not automatically a coupling for the union. Detecting the
disjoint case cheaply, and finding a principled combination when footprints
overlap, is the open problem. Its payoff would be direct. The traversal could fuse
the safe sets of neighbouring cells and prune the subset space far more
aggressively than certifying each in isolation allows.

\section{Symmetry-Orbit Space Exploration} % automorphisms / relabeling
\label{sec:fw-symmetry}

The subset-space traversal treats all wires alike, but masked gadgets are highly
symmetric. Permuting the share indices of an input, or relabelling interchangeable
masks, carries one probe to another with identical secret-dependence. Such a
relabelling is a circuit automorphism, and the automorphisms partition the probes
into orbits on which the verdict is constant, so verifying one representative per
orbit would certify the entire orbit and shrink the space the traversal must
explore by roughly the size of the symmetry group. Cutting a search by exploiting
its symmetry this way is a known technique in model checking
\parencite{clarke1996exploiting}. Carrying it into our setting, by detecting a
gadget's automorphisms and quotienting the Vandermonde recursion of
Section~\ref{sec:vandermonde} by them, is a promising and, to our knowledge,
largely unexplored direction for language-based probing-security verification.

\section{Scaling and Compositional Verification}
\label{sec:fw-composition}

We have previously mentioned that probing security does
not compose, so the approach does not, on its own, scale to a full cipher assembled
from gadgets. The composability notions surveyed in
Section~\ref{sec:related-composition} are the established works in this line. Strong
non-interference \parencite{barthe2016strong} and probe-isolating non-interference
\parencite{cassiers2020trivially} let one certify small gadgets once and wire them
together with the order preserved, while tight approaches
\parencite{belaid2018tight} track exactly the refreshing this needs. Our engine
could serve as the per-gadget oracle inside such a scheme, discharging each
gadget's stronger non-interference obligation while a compositional layer handles
the assembly. Adapting the work to include such compositional methods is work that remains
to be explored.

\section{Glitches and Practical Aspects}
\label{sec:fw-glitches}

Our model is value-based (Chapter~\ref{chap:preliminaries}). So, a probed wire leaks
its logical value and nothing more. Any hardware designer knows that reality is a lot less forgiving. A gate
glitches on unsettled inputs and a register transition leaks a function of two
successive values, and either can reopen a leak that a value-based check certifies.
The glitch-extended probing model, formalised together with composition in
the robust probing model \parencite{faust2018robust}, captures these effects, and
some tools verify masking under them \parencite{bloem2018formal,
10.1007/978-3-030-29959-0_15}, with recent work spanning both Boolean and
arithmetic masking across hardware and software \parencite{gigerl2023formal}.
Lifting our checker to a glitch-extended model would mean charging a probe on a
glitching gate with all of that gate's stable inputs, enlarging each observation
before the engine ever sees it. The coupling machinery would carry over easily as well;
what changes is the front end that forms the probes. Validating the result on real
netlists is the practical work that would turn the tool from a checker of simple circuits
into a verifier of practical implementations.
